% ===================== Chapter 6 =====================
\chapter{Functions}
\label{ch:functions}

We have used \code{main} since the first page, and a few helper functions along
the way. Now we study functions properly, because they are the single most
important tool for organising a program. A \textbf{function} is a named,
reusable piece of code that takes some inputs, does some work, and (often) hands
back a result. Functions let you write a piece of logic once, give it a name,
and call it as many times as you like.

\section{Defining and calling a function}

You have already seen the shape. A function definition is the keyword
\code{func}, a name, a parenthesised list of parameters, an optional return
type, a colon, the body, and \code{end}:

\begin{salam}
func square(n i32) i32:
    ret n * n
end

func main:
    println square(5)
    println square(9)
end
\end{salam}

\begin{progout}
25
81
\end{progout}

\code{square} takes one \textbf{parameter}, \code{n} of type \code{i32}, and its
\textbf{return type} (written after the parameter list) is also \code{i32}. The
body computes \code{n * n} and hands it back with \code{ret}. In \code{main} we
\textbf{call} the function by writing its name and an \textbf{argument} in
parentheses: \code{square(5)}. The argument \code{5} becomes the parameter
\code{n} for that call.

\begin{note}
\emph{Parameter} and \emph{argument} are two ends of the same connection. The
parameter is the name in the definition (\code{n}); the argument is the value you
pass at the call site (\code{5}). The argument is copied into the parameter when
the call begins.
\end{note}

\section{Returning a value}

The \code{ret} keyword does two things at once: it ends the function immediately
and sends a value back to the caller. The type of that value must match the
function's declared return type.

A function can have several \code{ret} statements useful when different cases
produce different answers. The first one reached wins:

\begin{salam}
func abs(n i32) i32:
    if n < 0:
        ret -n
    end
    ret n
end

func main:
    println abs(-7), abs(7)
end
\end{salam}

\begin{progout}
7 7
\end{progout}

If \code{n} is negative, the first \code{ret} fires and the function is done; the
final \code{ret n} is reached only for non-negative inputs. Using an early
\code{ret} to dispatch a special case, then letting the main path continue
un-indented, is a clean and common style.

\section{Functions that return nothing}

Some functions \emph{do} something rather than \emph{compute} something they
print, they update, they save a file. Such a function has no return type: you
simply omit it after the parameter list. Its (implicit) result type is
\code{void}, the ``no value'' type from Chapter~2.

\begin{salam}
func greet(name str):
    println "Hello,", name
    println "Welcome to Salam."
end

func main:
    greet("Sara")
    greet("Ali")
end
\end{salam}

\begin{progout}
Hello, Sara
Welcome to Salam.
Hello, Ali
Welcome to Salam.
\end{progout}

Inside a \code{void} function you may still use a bare \code{ret} (with no value)
to return early for instance to bail out when an input is invalid but you
never return a value from one.

\section{Multiple parameters}

A function may take any number of parameters, separated by commas. Each needs its
own type:

\begin{salam}
func add(a i32, b i32) i32:
    ret a + b
end

func max3(a i32, b i32, c i32) i32:
    mut m i32 = a
    if b > m: m = b end
    if c > m: m = c end
    ret m
end

func main:
    println add(3, 4)
    println max3(7, 2, 9)
end
\end{salam}

\begin{progout}
7
9
\end{progout}

Arguments are matched to parameters \emph{by position}: in \code{add(3, 4)},
\code{3} becomes \code{a} and \code{4} becomes \code{b}.

\begin{warn}
The current version of Salam does not support \emph{default} parameter values, so
every parameter must be supplied at every call. Calling \code{add(3)} is an
error. (To offer a ``default,'' write a second overload see below.)
\end{warn}

\section{Arguments are passed by value}

When you pass an argument, the function receives a \emph{copy} of it. Changing a
parameter inside the function therefore does \emph{not} change the caller's
variable:

\begin{salam}
func tryToChange(x i32):
    mut local i32 = x
    local = 999            // only the copy changes
end

func main:
    n i32 = 5
    tryToChange(n)
    println n              // still 5
end
\end{salam}

\begin{progout}
5
\end{progout}

This ``pass by value'' rule makes functions easy to reason about: a function
cannot secretly reach back and alter the variables you passed in. When you want a
function to produce a new value, the clean approach is to \emph{return} it and let
the caller decide what to do:

\begin{salam}
func incremented(x i32) i32:
    ret x + 1
end

func main:
    mut n i32 = 5
    n = incremented(n)     // the caller chooses to update n
    println n
end
\end{salam}

\begin{progout}
6
\end{progout}

\begin{note}
Salam deliberately has no ``pass by reference'' for ordinary variables you
cannot take the address of a local with \code{\&} and hand it to a function. This
keeps the flow of data visible: values go \emph{in} as arguments and come
\emph{out} as return values. (In Chapter~9 you will see that a struct's own
methods can update the struct they belong to; that is the sanctioned way to
modify data in place.)
\end{note}

\section{Recursion}

A function is allowed to call \emph{itself}. This is called \textbf{recursion},
and it is the natural way to express anything defined in terms of a smaller
version of itself. The classic example is the factorial, where
$n! = n \times (n-1)!$:

\begin{salam}
func fact(n i32) i32:
    if n <= 1:
        ret 1
    end
    ret n * fact(n - 1)
end

func main:
    println fact(5)
end
\end{salam}

\begin{progout}
120
\end{progout}

Every recursive function needs two things to be correct: a \textbf{base case}
that stops the recursion (here, \code{n <= 1} returns \code{1} without calling
again), and a \textbf{recursive step} that moves \emph{toward} the base case
(here, calling \code{fact(n - 1)} with a smaller argument). Trace \code{fact(5)}:
it becomes \code{5 * fact(4)}, which becomes \code{5 * 4 * fact(3)}, and so on
down to \code{fact(1)}, which is \code{1} then the multiplications unwind:
$5\times4\times3\times2\times1 = 120$.

\begin{warn}
Forget the base case or fail to shrink the argument and a recursive
function calls itself forever, until the program runs out of stack space and
crashes. Always ask: ``what stops this?'' and ``does each call get closer to
stopping?''
\end{warn}

\begin{deep}[In Depth: recursion versus iteration]
Anything you can compute with recursion you can also compute with a loop, and
vice versa. Recursion often expresses a problem more clearly (tree walks, divide
and conquer), while a loop avoids the small cost of repeated function calls. For
factorials a loop is perfectly natural:
\begin{salam}
func fact(n i32) i32:
    mut result i32 = 1
    for mut i = 2; i <= n; i = i + 1:
        result = result * i
    end
    ret result
end
\end{salam}
Choose whichever makes the intent clearest; both are idiomatic Salam.
\end{deep}

\section{Definition order does not matter}

You may call a function that is defined \emph{later} in the file. Salam reads all
the top-level definitions before it runs anything, so \code{main} can freely call
helpers written below it:

\begin{salam}
func main:
    println later()
end

func later() str:
    ret "defined after main"
end
\end{salam}

\begin{progout}
defined after main
\end{progout}

This means you never need ``forward declarations,'' and you are free to put
\code{main} at the top of the file where a reader expects to find it.

\section{Function overloading}

Two functions may share a name as long as their parameters differ in number or
in type. This is called \textbf{overloading}, and Salam picks the right one based
on the arguments at each call:

\begin{salam}
func add(a i32, b i32) i32:
    ret a + b
end

func add(a f64, b f64) f64:
    ret a + b
end

func main:
    println add(2, 3)        // calls the i32 version
    println add(1.5, 2.5)    // calls the f64 version
end
\end{salam}

\begin{progout}
5
4
\end{progout}

The two \code{add}s are genuinely different functions that happen to share a name.
Overloading lets a family of related operations present one consistent name
instead of \code{addInt}, \code{addFloat}, and so on. (The second result prints as
\code{4} because it is the floating value \code{4.0}.)

\section{Use every parameter}

Salam is tidy about parameters: if you declare one and never use it, the compiler
reports an error rather than letting the dead weight slip by. If you genuinely
need a parameter you do not use perhaps to match an expected shape prefix
its name with an underscore to say ``unused on purpose'':

\begin{salam}
func ignore_second(a i32, _b i32) i32:
    ret a
end
\end{salam}

\begin{tip}
This rule nudges you toward honest signatures: a function's parameter list should
list exactly what it needs. When the compiler flags an unused parameter, it is
often pointing at a leftover from an earlier version of the code that you can
simply delete.
\end{tip}

\section{Designing good functions}

Functions are how you tame complexity. A few guidelines that pay off:

\begin{itemize}[leftmargin=1.4em]
  \item \textbf{One job each.} A function should do a single, nameable thing. If
        you struggle to name it, it is probably doing too much.
  \item \textbf{Short is good.} If a function grows past a screenful, look for a
        piece you can lift out into its own well-named helper.
  \item \textbf{Name for the result.} \code{is\_leap(year)} and
        \code{average(scores)} tell the reader what comes back. A good name turns
        a call site into a readable sentence.
  \item \textbf{Prefer returning over reaching out.} Compute a value and return
        it, rather than relying on shared mutable state.
\end{itemize}

\section{Chapter summary}

\begin{itemize}[leftmargin=1.4em]
  \item A function is \code{func name(params) RetType:} \dots\ \code{end};
        \code{ret} returns a value and ends the function.
  \item Omit the return type for a function that returns nothing (\code{void});
        a bare \code{ret} returns early from one.
  \item Parameters are matched by position and passed \emph{by value} (the
        function gets a copy). There are no default parameters and no reference
        parameters.
  \item A function may call itself (recursion); give it a base case and make each
        call smaller.
  \item Definition order is irrelevant you can call functions defined later.
  \item Functions may be \emph{overloaded} by differing parameters.
  \item Every declared parameter must be used, or prefixed with \code{\_}.
\end{itemize}

\section{Exercises}

\exercise Write a function \code{cube(n i32) i32} that returns $n^3$ using
multiplication, and print the cubes of 1 through 5 with a loop.

\exercise Write a \code{void} function \code{banner(text str)} that prints the
text surrounded by a line of dashes above and below it.

\exercise Write \code{min2(a i32, b i32) i32} and use it to build
\code{min3(a, b, c)} by calling \code{min2} twice.

\exercise Write a recursive function \code{sum\_to(n i32) i32} that returns
$1 + 2 + \dots + n$. Identify its base case and recursive step in comments.

\exercise Overload a function \code{describe} so that \code{describe(42)} prints
``an integer'' and \code{describe(3.14)} prints ``a float''. Test both.

\exercise Rewrite the recursive \code{sum\_to} from a previous exercise as a loop,
and confirm both give the same answers.
